\documentclass[12pt]{article}

\usepackage{amsmath}
\usepackage{amssymb}
\usepackage{graphicx}
\graphicspath{{figures/}{figures/fft/}{figures/cosmo_ic/}}
\usepackage{hyperref}
\usepackage[margin=1in]{geometry}
\usepackage{natbib}
\usepackage{listings}
\usepackage{xcolor}
\usepackage{booktabs}
\usepackage{tcolorbox}
\tcbuselibrary{breakable,skins}

% --- code listing style ---
\lstset{
  basicstyle=\ttfamily\small,
  keywordstyle=\color{blue},
  commentstyle=\color{gray},
  stringstyle=\color{teal},
  breaklines=true,
  frame=single,
  numbers=left,
  numberstyle=\tiny\color{gray}
}

\newcommand{\bk}{\mathbf{k}}
\newcommand{\bx}{\mathbf{x}}
\newcommand{\bn}{\mathbf{n}}
\newcommand{\ba}{\mathbf{a}}
\newcommand{\by}{\mathbf{y}}
\newcommand{\bq}{\mathbf{q}}
\newcommand{\br}{\mathbf{r}}
\newcommand{\bz}{\mathbf{z}}
\newcommand{\bu}{\mathbf{u}}
\newcommand{\bff}{\mathbf{f}}
\newcommand{\avg}[1]{\left\langle #1 \right\rangle}
\newcommand{\code}[1]{\texttt{#1}}

\title{Statistical Foundations for Cosmological Fields:\\
       Random Fields, Correlation Functions, Power Spectra,\\
       and the Fourier Machinery that Links Them}
\author{}
\date{\today}

\begin{document}

\maketitle

\begin{tcolorbox}[colback=yellow!10!white, colframe=orange!75!black,
                  title={\bfseries Disclaimer}]
These notes were compiled by an AI assistant (Claude, Anthropic) and have
\textbf{not been reviewed by a human domain expert}.  They may contain
errors in derivations, physical reasoning, or numerical estimates.
Cross-check key equations against the primary literature before use.
\end{tcolorbox}
\vspace{0.5em}

\begin{tcolorbox}[colback=blue!5!white, colframe=blue!60!black,
                  title={\bfseries What this note is for}]
This is the ``chapter zero'' for the rest of the cosmology notes in this
directory.  Before one can talk about the initial conditions of a
simulated universe, one needs a language for describing a quantity ---
the density --- that is different in every realisation and yet
statistically predictable.  That language is the theory of \emph{random
fields}, and its two central summary statistics are the
\emph{correlation function} and the \emph{power spectrum}, which are two
views of the same object connected by a \emph{Fourier transform}.  This
note builds all of that from scratch, with full derivations and physical
interpretation, assuming only calculus and a first course in
probability.  Read it before \code{cosmo\_ic.tex} (cosmological initial
conditions), the numerical Fourier transform, and
\code{ic\_sampling.tex} (sampling methods).  Wherever a topic is
developed more deeply in one of those notes, a pointer is given.
\end{tcolorbox}

\tableofcontents
\newpage

%======================================================================
\section{Why cosmology needs a theory of random fields}
\label{sec:why}

\subsection{One universe, many possible universes}

Look at a map of the galaxies, or of the temperature of the microwave
background, or of the dark-matter density in a simulation.  What you see
is a specific, bumpy, irregular pattern.  If you asked ``what is the
density of matter at this exact point in space?'' the honest answer is
that the theory does not predict it and never will.  The early universe
laid down small fluctuations through a quantum-mechanical process during
inflation, and those fluctuations are, as far as we can tell, drawn at
random.  We live in one particular outcome of that random process.

This is not a failure of the theory; it is the nature of the object.
The right analogy is a rolled die or a noisy voltage on a wire.  You
cannot predict the next roll or the exact voltage at the next instant,
but you can predict the \emph{statistics}: the probability of each face,
the average voltage, the size of the fluctuations, how the voltage now
is related to the voltage a microsecond later.  Cosmology is in exactly
this situation, with one extra twist that we will have to confront: we
get to observe only \emph{one} roll of the cosmic die --- a single
universe --- and yet we want to test statistical predictions.  Making
that possible is the job of \emph{ergodicity} (\S\ref{sec:ergodic}).

\subsection{The density contrast as the basic field}

Let $\rho(\bx)$ be the mass density at comoving position $\bx$, and let
$\bar\rho$ be its spatial average.  Absolute densities are awkward
because the mean changes as the universe expands, so we work with the
dimensionless \emph{density contrast}
\begin{equation}
  \delta(\bx) \;\equiv\; \frac{\rho(\bx) - \bar\rho}{\bar\rho}\,.
  \label{eq:delta-def}
\end{equation}
By construction $\delta \ge -1$ (density cannot go negative), and its
spatial average is zero: $\avg{\delta} = 0$.  Where matter has piled up,
$\delta > 0$; in a void, $\delta < 0$.  In the early universe $|\delta|
\ll 1$ everywhere, which is what makes the linear theory of the
companion notes so powerful.

A quantity like $\delta(\bx)$ --- a number attached to every point of
space, whose value at each point is a random variable --- is called a
\emph{random field}.  Everything in this note is a statement about how
to describe such an object, how to compress it into a few meaningful
functions, and how to build a computer realisation of one.

The same formalism, essentially unchanged, describes the peculiar
velocity field, the gravitational potential, the temperature of the
microwave background on the sky, weak-lensing shear, and the pressure
of a turbulent fluid.  Nothing below is special to cosmology except the
choice of examples.

%======================================================================
\section{Ensembles, averages, and the symmetries we assume}
\label{sec:ensemble}

\subsection{The ensemble average}
\label{sec:ensemble-avg}

Imagine not one universe but an infinite collection of them, all
generated by the same physical process with the same statistical rules,
differing only in the particular random numbers that came up.  This
imagined collection is the \emph{ensemble}.  The \emph{ensemble average}
of a quantity, written $\avg{\,\cdot\,}$, is the average over all members
of the ensemble at a \emph{fixed} spatial point.  For example
$\avg{\delta(\bx)}$ means: fix the point $\bx$, look at its value of
$\delta$ in every universe of the ensemble, and average those numbers.

This is a theoretical construct --- we have only one universe --- but it
is the correct object for stating predictions, because the physics
(inflation, gravity) predicts the \emph{probability rules} of the
ensemble, not the outcome in any one member.  All the statistics below
(mean, correlation function, power spectrum) are ensemble averages.
Section~\ref{sec:ergodic} explains how we nonetheless estimate them
from our single realisation.

\subsection{The mean and why we can set it to zero}

The first statistic is the mean, $\avg{\delta(\bx)}$.  For the density
contrast it is zero everywhere by the definition
(\ref{eq:delta-def}).  More generally, whenever the mean is a known
deterministic function we can subtract it and work with the
fluctuation $\delta - \avg{\delta}$, which has zero mean.  So without
loss of generality we take
\begin{equation}
  \avg{\delta(\bx)} = 0 \qquad \text{for all } \bx .
\end{equation}
A zero-mean field carries \emph{no} information in its first moment; all
the structure is in the second moment and higher, which is where we go
next.

\subsection{Statistical homogeneity}
\label{sec:homogeneity}

The cosmological principle asserts that the universe has no special
places: the statistical rules are the same everywhere.  Translated into
the language of the ensemble, this is \emph{statistical homogeneity} ---
the requirement that every ensemble average is unchanged if we shift all
positions by a common displacement $\ba$:
\begin{equation}
  \avg{\delta(\bx_1)\,\delta(\bx_2)\cdots}
  \;=\;
  \avg{\delta(\bx_1+\ba)\,\delta(\bx_2+\ba)\cdots}
  \qquad\text{for every }\ba .
  \label{eq:homog}
\end{equation}
Note carefully what homogeneity does \emph{not} say.  It does not say
$\delta$ is the same everywhere --- that would be a constant field with
no structure.  It says the \emph{statistics} are the same everywhere:
a void near here is as likely as a void of the same size near there.
Any one realisation is still as bumpy as ever.

The immediate consequence, applied to the two-point average, is that
$\avg{\delta(\bx)\,\delta(\bx')}$ can depend only on the separation
$\br = \bx' - \bx$, not on where the pair sits.  This is what makes the
correlation function a function of one vector argument instead of two,
and it is the single most-used assumption in this note.

\subsection{Statistical isotropy}
\label{sec:isotropy}

A stronger symmetry is \emph{statistical isotropy}: no special
directions either.  Formally, every ensemble average is invariant under
a common rotation $R$ of all positions,
\begin{equation}
  \avg{\delta(R\bx_1)\,\delta(R\bx_2)\cdots}
  = \avg{\delta(\bx_1)\,\delta(\bx_2)\cdots}.
\end{equation}
Combined with homogeneity, isotropy forces the two-point function to
depend only on the \emph{magnitude} of the separation, $r = |\br|$, not
its direction.  So under the two symmetries together the entire
two-point content of the field is a single function of a single scalar,
$\xi(r)$.  That enormous compression --- from ``a value at every pair of
points'' to ``one curve'' --- is why these statistics are so useful.

Isotropy can be broken in real observations (redshift-space distortions
pick out the line of sight; a bulk flow picks out a direction), and then
one keeps the direction dependence.  For the initial conditions we care
about here, isotropy holds to excellent accuracy and we use it freely.

\subsection{Ergodicity: measuring an ensemble from one realisation}
\label{sec:ergodic}

Here is the twist promised in \S\ref{sec:why}.  Every statistic defined
so far is an ensemble average --- an average over universes we do not
have.  We have one universe.  How can we possibly measure
$\avg{\delta\,\delta}$?

The rescue is \emph{ergodicity}: the assumption that widely separated
regions of a single realisation are statistically independent, so that
they behave like independent draws from the ensemble.  If that holds,
then averaging a quantity over the volume of one large realisation gives
the same answer as averaging over the ensemble at a point:
\begin{equation}
  \underbrace{\avg{\,\cdot\,}}_{\text{average over universes}}
  \;=\;
  \underbrace{\frac{1}{V}\int_V \mathrm d^d x\;(\,\cdot\,)}_{\text{average over one universe}}
  \qquad (V \to \infty).
  \label{eq:ergodic}
\end{equation}
Physically the condition is that the correlation function falls to zero
at large separation (\S\ref{sec:xi-props}): once two patches are farther
apart than the correlation length, they know nothing about each other
and count as independent samples.  A box many correlation lengths across
therefore contains many effectively-independent sub-volumes, and the
spatial average over it converges to the ensemble average.  The residual
error from having only a finite number of independent patches is
\emph{cosmic variance}, and it is the fundamental noise floor on any
measurement of $\xi$ or $P$ from a finite volume (\S\ref{sec:cosmicvar}).

Every measured correlation function and power spectrum in cosmology ---
and every one computed from a simulation box --- relies on
Eq.~(\ref{eq:ergodic}).  Keep it in mind: it is the bridge between the
theory (ensemble language) and the data (one box).

%======================================================================
\section{The two-point correlation function}
\label{sec:xi}

\subsection{Definition and physical meaning}
\label{sec:xi-def}

The \emph{two-point correlation function} is the ensemble average of the
product of the field at two points.  Using homogeneity
(\S\ref{sec:homogeneity}) to reduce it to a function of the separation
$\br = \bx' - \bx$,
\begin{equation}
  \boxed{\;\xi(\br) \;\equiv\; \avg{\delta(\bx)\,\delta(\bx+\br)}\;}
  \label{eq:xi-def-box}
\end{equation}
with the right-hand side independent of $\bx$.  Under isotropy it
reduces further to $\xi(r)$ with $r = |\br|$.

What does it mean?  Suppose you stand at a point where the density is
above average ($\delta(\bx) > 0$).  The correlation function tells you
whether a point a distance $r$ away tends to be above or below average
too.  If $\xi(r) > 0$, the two points are \emph{positively correlated}:
an overdensity here makes an overdensity there more likely --- matter
clusters.  If $\xi(r) < 0$, they are anti-correlated: an overdensity
here is compensated by an underdensity there.  If $\xi(r) = 0$, the two
points are statistically independent.

There is a clean interpretation in terms of counting pairs.  For a
point process (galaxies, particles) of mean number density $\bar n$, the
probability of finding a particle in a small volume $\mathrm dV_1$ and
another in $\mathrm dV_2$ separated by $r$ is
\begin{equation}
  \mathrm dP_{12} = \bar n^2\,\bigl[1 + \xi(r)\bigr]\,\mathrm dV_1\,\mathrm dV_2 .
  \label{eq:xi-pair}
\end{equation}
The ``1'' is what you would get if particles were sprinkled completely
at random (a Poisson process); the $\xi(r)$ is the \emph{excess}
probability of finding a pair at separation $r$ over and above random.
This is why $\xi$ is often called the excess clustering: $\xi(r) = 0.5$
means pairs at separation $r$ are $50\%$ more common than chance.

\subsection{Basic properties}
\label{sec:xi-props}

Three properties follow directly from the definition
(\ref{eq:xi-def-box}) and are worth committing to memory.

\paragraph{(i) At zero separation, $\xi$ is the variance.}
Setting $\br = 0$,
\begin{equation}
  \xi(0) = \avg{\delta(\bx)^2} = \sigma^2 ,
  \label{eq:xi-zero}
\end{equation}
the variance of the field.  It measures the typical squared amplitude of
the fluctuations: $\sigma = \sqrt{\xi(0)}$ is the root-mean-square
density contrast.

\paragraph{(ii) It is symmetric.}
Because the two points enter symmetrically and the field is real,
$\xi(\br) = \xi(-\br)$; under isotropy this is automatic since $\xi$
depends only on $r$.

\paragraph{(iii) It decays at large separation.}
Distant points are physically unrelated, so $\xi(r) \to 0$ as $r \to
\infty$.  The scale over which it decays is the \emph{correlation
length}.  This decay is exactly the condition that makes ergodicity
(\S\ref{sec:ergodic}) work: beyond the correlation length, patches are
independent.

\subsection{The correlation function is a covariance}
\label{sec:xi-covariance}

It pays to see $\xi$ through the lens of ordinary statistics.  For two
random variables $X$ and $Y$ with zero mean, the \emph{covariance} is
$\mathrm{Cov}(X, Y) = \avg{XY}$.  Comparing with
Eq.~(\ref{eq:xi-def-box}), the correlation function is nothing but the
covariance between the field values at two points:
\begin{equation}
  \xi(\bx - \bx') = \mathrm{Cov}\bigl(\delta(\bx),\,\delta(\bx')\bigr).
\end{equation}
If we discretise space into $N$ cells and stack the field values into a
vector $\boldsymbol\delta = (\delta_1, \dots, \delta_N)^\top$, the full
second-moment information is the \emph{covariance matrix}
\begin{equation}
  C_{ij} \;=\; \avg{\delta_i\,\delta_j} \;=\; \xi(\bx_i - \bx_j).
  \label{eq:covmatrix}
\end{equation}
So ``the correlation function'' and ``the covariance matrix of the
discretised field'' are the same object.  This viewpoint --- $\xi$ as a
matrix --- is the key to two things later: it explains why the power
spectrum is the natural diagonal basis (\S\ref{sec:cov-diag}), and it is
the starting point for actually generating a field on a computer
(\S\ref{sec:generate}).  Homogeneity, in this language, says that $C_{ij}$
depends only on the offset $i - j$; a matrix with that property is called
\emph{Toeplitz}, or, on a periodic grid, \emph{circulant}, and circulant
matrices have a magical diagonalisation that we will exploit.

%======================================================================
\section{Interlude: the Fourier transform and its conventions}
\label{sec:fourier}

The power spectrum is the correlation function seen in Fourier space, so
we need the Fourier transform first.  This section fixes conventions and
states the few facts we use; a standard reference on the discrete Fourier transform develops
the numerical (discrete, fast) transform in full, and Appendix~A of
\code{cosmo\_ic.tex} carries out the careful bookkeeping of factors that
we here quote.

\subsection{The continuous Fourier transform}
\label{sec:cft}

Any well-behaved function $f(\bx)$ on $\mathbb R^d$ can be written as a
superposition of plane waves $e^{i\bk\cdot\bx}$.  We use the symmetric
``physics'' convention, in which all the $2\pi$'s sit in the inverse
transform:
\begin{align}
  \hat f(\bk) &= \int \mathrm d^d x\; f(\bx)\,e^{-i\bk\cdot\bx},
  \label{eq:ft}\\
  f(\bx) &= \int \frac{\mathrm d^d k}{(2\pi)^d}\; \hat f(\bk)\,e^{+i\bk\cdot\bx}.
  \label{eq:ift}
\end{align}
The wavevector $\bk$ has units of inverse length; its magnitude $k =
|\bk|$ is the wavenumber and $2\pi/k$ is the wavelength of that plane
wave.  A large $k$ is a short-wavelength, small-scale feature; a small
$k$ is a long-wavelength, large-scale feature.  This dictionary ---
\emph{large $k$ means small scale} --- is worth internalising because
cosmological intuition is usually phrased in terms of scales but the
mathematics lives in $k$.

The one identity we lean on repeatedly is the integral representation of
the Dirac delta,
\begin{equation}
  \int \mathrm d^d x\; e^{-i\bq\cdot\bx} = (2\pi)^d\,\delta^D(\bq),
  \label{eq:delta-rep}
\end{equation}
which is just Eqs.~(\ref{eq:ft})--(\ref{eq:ift}) applied to $f = 1$.  It
says: integrate a pure plane wave over all space and you get zero unless
the wave is the constant one ($\bq = 0$), in which case you get an
(infinite) spike.  Distinct plane waves are orthogonal.

\subsection{Fourier series in a finite box}
\label{sec:fseries}

Simulations and real surveys live in a finite volume, not on all of
$\mathbb R^d$.  On a cubic box of side $L$ and volume $V = L^d$ with
periodic boundaries, only wavevectors that fit an integer number of
wavelengths in the box are allowed:
\begin{equation}
  \bk = \frac{2\pi}{L}\,\bn, \qquad \bn \in \mathbb Z^d .
\end{equation}
The smallest nonzero wavenumber is the \emph{fundamental mode}
$k_{\rm fund} = 2\pi/L$: no fluctuation longer than the box can be
represented.  The transform becomes a sum over these discrete modes,
\begin{equation}
  \delta(\bx) = \frac{1}{V}\sum_{\bk} \hat\delta_{\bk}\,e^{i\bk\cdot\bx},
  \qquad
  \hat\delta_{\bk} = \int_V \mathrm d^d x\; \delta(\bx)\,e^{-i\bk\cdot\bx},
  \label{eq:fseries}
\end{equation}
and the continuum rule $\int \mathrm d^d k/(2\pi)^d \to V^{-1}\sum_\bk$
converts between the two pictures.  The box turns the continuous Dirac
delta of Eq.~(\ref{eq:delta-rep}) into a Kronecker delta,
$\int_V \mathrm d^d x\, e^{-i(\bk-\bk')\cdot\bx} = V\,\delta^K_{\bk\bk'}$.

\subsection{Reality forces Hermitian symmetry}
\label{sec:hermitian}

Our fields are real: $\delta(\bx)$ is a physical density contrast, not a
complex number.  Take the complex conjugate of Eq.~(\ref{eq:ft}); since
$\delta$ is real, $\delta^* = \delta$, and the only change on the
right-hand side is $e^{-i\bk\cdot\bx} \to e^{+i\bk\cdot\bx}$, which is
the transform evaluated at $-\bk$.  Hence
\begin{equation}
  \boxed{\;\hat\delta(-\bk) = \hat\delta^*(\bk)\;}
  \label{eq:hermitian}
\end{equation}
the \emph{Hermitian symmetry} of a real field's transform.  It says the
mode at $-\bk$ is not independent --- it is the complex conjugate of the
mode at $+\bk$.  Only half of $k$-space carries independent information;
the other half is determined.  This has a concrete consequence when we
generate fields (\S\ref{sec:generate}): if you draw the Fourier
coefficients by hand you must impose Eq.~(\ref{eq:hermitian}), whereas if
you build the field from real-space white noise it comes out
automatically.  A handful of special modes (the $\bk = 0$ mode, and the
Nyquist modes on an even grid) are their own conjugates and must
therefore be purely real; the discrete-transform literature works
through those corner cases.

%======================================================================
\section{The power spectrum}
\label{sec:pk}

\subsection{From the correlation of modes to $P(k)$}
\label{sec:pk-derive}

We now compute the two-point statistics in Fourier space and discover
that homogeneity makes different modes uncorrelated.  Take two Fourier
coefficients and average their product, using Eq.~(\ref{eq:ft}) for
each:
\begin{align}
  \avg{\hat\delta(\bk)\,\hat\delta^*(\bk')}
  &= \int \mathrm d^d x\,\mathrm d^d x'\;
     \avg{\delta(\bx)\,\delta(\bx')}\;
     e^{-i\bk\cdot\bx + i\bk'\cdot\bx'} \notag\\
  &= \int \mathrm d^d x\,\mathrm d^d x'\;
     \xi(\bx' - \bx)\;
     e^{-i\bk\cdot\bx + i\bk'\cdot\bx'},
  \label{eq:pk-step1}
\end{align}
where the second line used homogeneity to replace the correlation with
$\xi$ of the separation.  Now change variables from $(\bx, \bx')$ to
$(\bx, \br)$ with $\br = \bx' - \bx$ --- a pure shift, so the Jacobian is
one and $\mathrm d^d x\,\mathrm d^d x' = \mathrm d^d x\,\mathrm d^d r$.
The exponent becomes
$-i\bk\cdot\bx + i\bk'\cdot(\bx + \br) = -i(\bk - \bk')\cdot\bx +
i\bk'\cdot\br$, and the double integral factorises:
\begin{align}
  \avg{\hat\delta(\bk)\,\hat\delta^*(\bk')}
  &= \underbrace{\int \mathrm d^d x\; e^{-i(\bk-\bk')\cdot\bx}}_{(2\pi)^d\,\delta^D(\bk-\bk')}
     \;\;\underbrace{\int \mathrm d^d r\; \xi(\br)\,e^{i\bk'\cdot\br}}_{\text{a function of }\bk'} \notag\\
  &= (2\pi)^d\,\delta^D(\bk - \bk')\;P(\bk),
  \label{eq:pk-def}
\end{align}
where the first integral is the delta identity (\ref{eq:delta-rep}) and
we \emph{defined} the quantity multiplying it as the power spectrum,
\begin{equation}
  \boxed{\;P(\bk) \;\equiv\; \int \mathrm d^d r\; \xi(\br)\,e^{-i\bk\cdot\br}\;}
  \label{eq:wiener-khinchin}
\end{equation}
(the sign of the exponent flipped because the delta forced $\bk' = \bk$).
Two results have dropped out at once.

\paragraph{Result 1: distinct modes are uncorrelated.}
The factor $\delta^D(\bk - \bk')$ in Eq.~(\ref{eq:pk-def}) means
$\avg{\hat\delta(\bk)\,\hat\delta^*(\bk')} = 0$ unless $\bk = \bk'$.
Homogeneity \emph{diagonalises} the field in Fourier space: different
Fourier modes are statistically independent.  This is the deep reason
Fourier space is the natural home for a homogeneous random field, and
it is the wave-space echo of the circulant-matrix remark at the end of
\S\ref{sec:xi-covariance}.

\paragraph{Result 2: the power spectrum is the Fourier transform of the
correlation function.}  Equation~(\ref{eq:wiener-khinchin}) is the
\emph{Wiener--Khinchin theorem}: $P$ and $\xi$ are a Fourier-transform
pair,
\begin{equation}
  P(\bk) = \int \mathrm d^d r\;\xi(\br)\,e^{-i\bk\cdot\br},
  \qquad
  \xi(\br) = \int \frac{\mathrm d^d k}{(2\pi)^d}\;P(\bk)\,e^{+i\bk\cdot\br}.
  \label{eq:P-xi-pair}
\end{equation}
They are the same information in two languages.  The correlation
function answers ``how related are two points a distance $r$ apart?'';
the power spectrum answers ``how much fluctuation power is there on
scale $2\pi/k$?''  Neither contains anything the other lacks; which you
use is a matter of convenience, and often of which one the measurement
apparatus hands you more directly.

\paragraph{Three properties of $P(k)$.}  These follow immediately and are
worth stating.  (i) \emph{$P$ is real.}  Because $\xi(\br)$ is real and
even (\S\ref{sec:xi-props}), its transform
$\int \xi(\br)\,e^{-i\bk\cdot\br}\,\mathrm d^d r$ has an imaginary part
$\propto \int \xi(\br)\sin(\bk\cdot\br)\,\mathrm d^d r$ that integrates to
zero by the oddness of the sine against the even $\xi$; only the cosine
(real) part survives.  (ii) \emph{$P$ is even}, $P(-\bk) = P(\bk)$, for
the same reason, and under isotropy it depends only on $k = |\bk|$.
(iii) \emph{$P$ is non-negative}, $P(\bk) \ge 0$.  This is not obvious
from Eq.~(\ref{eq:wiener-khinchin}) --- it is forced by the
variance-per-mode form (\ref{eq:pk-var}) below, since a variance cannot
be negative.  The non-negativity is a genuine constraint: not every even
function is a legal correlation function; only those whose Fourier
transform is everywhere non-negative are (this is \emph{Bochner's
theorem}, the statement that $\xi$ must be a positive-definite function,
equivalently that the covariance matrix of \S\ref{sec:xi-covariance} is
positive-semidefinite).

\subsection{The variance-per-mode form}
\label{sec:pk-variance}

On the finite box of \S\ref{sec:fseries} there is an even cleaner form,
derived directly rather than through the continuum delta.  The box
Fourier coefficient is $\hat\delta_\bk = \int_V \mathrm d^d x\,
\delta(\bx)\,e^{-i\bk\cdot\bx}$, so its mean-square is
\begin{align}
  \avg{|\hat\delta_\bk|^2}
  &= \int_V\!\!\int_V \mathrm d^d x\,\mathrm d^d x'\;
     \avg{\delta(\bx)\,\delta(\bx')}\;e^{-i\bk\cdot(\bx-\bx')} \notag\\
  &= \int_V \mathrm d^d x \int \mathrm d^d r\;\xi(\br)\,e^{-i\bk\cdot\br}
   \;=\; V\,P(\bk),
  \label{eq:pk-box-deriv}
\end{align}
where homogeneity again let the integrand depend only on
$\br = \bx - \bx'$, the $\br$-integral is $P(\bk)$ by
Eq.~(\ref{eq:wiener-khinchin}), and the leftover $\bx$-integral over the
box contributes the volume $V$ (the integrand no longer depends on
$\bx$).  Hence
\begin{equation}
  \boxed{\;P(\bk) = \frac{\avg{|\hat\delta_\bk|^2}}{V}\;}
  \label{eq:pk-var}
\end{equation}
The power spectrum is (up to the volume) the \emph{variance of the
Fourier amplitude of each mode}.  Modes with large $P(k)$ are the ones
that fluctuate strongly from one realisation to the next; modes with
small $P(k)$ are nearly quiet.  This is the form one actually measures:
square the Fourier amplitudes of the box, divide by the volume, and
average over modes of the same $|\bk|$ (\S\ref{sec:estimate-pk}).

\subsection{Physical meaning: power per scale and the variance budget}
\label{sec:pk-physical}

Set $\br = 0$ in the inverse relation of Eq.~(\ref{eq:P-xi-pair}) and use
$\xi(0) = \sigma^2$ from Eq.~(\ref{eq:xi-zero}):
\begin{equation}
  \sigma^2 = \xi(0) = \int \frac{\mathrm d^d k}{(2\pi)^d}\;P(\bk).
  \label{eq:sigma-integral}
\end{equation}
The total variance of the field is the integral of the power spectrum
over all modes.  Each mode contributes its little piece $P(\bk)\,\mathrm
d^d k/(2\pi)^d$ to the total fluctuation budget --- hence the name
``power spectrum'', in exact analogy with the power spectrum of a noisy
electrical signal, where the total power is the integral of the spectral
density over frequency.

In three dimensions, under isotropy, do the angular part of the integral
($\mathrm d^3 k = 4\pi k^2\,\mathrm d k$):
\begin{equation}
  \sigma^2 = \int_0^\infty \frac{4\pi k^2\,\mathrm d k}{(2\pi)^3}\,P(k)
           = \int_0^\infty \frac{k^3 P(k)}{2\pi^2}\;\frac{\mathrm d k}{k}
           = \int_0^\infty \Delta^2(k)\;\mathrm d\ln k,
  \label{eq:dimensionless}
\end{equation}
which defines the \emph{dimensionless power spectrum}
\begin{equation}
  \boxed{\;\Delta^2(k) \equiv \frac{k^3 P(k)}{2\pi^2}\;}
  \label{eq:Delta2}
\end{equation}
This combination is the variance contributed per logarithmic interval in
$k$ --- per octave of scale.  It is the single most useful diagnostic in
practice: if $\Delta^2(k) \approx 1$ at some $k$, then fluctuations on
that scale are of order unity and linear theory is breaking down there;
where $\Delta^2 \ll 1$ the field is still in the linear regime.  Plotting
$\Delta^2(k)$ rather than $P(k)$ puts equal-sized scales on equal
footing and makes the ``which scales have gone nonlinear'' question
immediate.

\subsection{Units and the $h$ bookkeeping}
From Eq.~(\ref{eq:wiener-khinchin}), $P(k)$ has units of
(length)$^d$ --- in 3D, volume.  Because $\xi$ is dimensionless for the
density contrast, $P(k)$ carries the $\mathrm d^3 r$.  Cosmologists
usually quote lengths in $h^{-1}\mathrm{Mpc}$, so $P(k)$ comes in
$(h^{-1}\mathrm{Mpc})^3$ and $k$ in $h\,\mathrm{Mpc}^{-1}$.  The
simulation code in this repository stores lengths in Mpc (no $h$); the
conversion is spelled out in \code{CLAUDE.md} and Appendix~B of
\code{cosmo\_ic.tex}.  Getting the factors of $h$ and $2\pi$ right is
tedious but mechanical; the safe practice is to fix one convention (we
use Eqs.~\ref{eq:ft}--\ref{eq:ift}) and carry it everywhere.

%======================================================================
\section{Covariance as the unifying object}
\label{sec:covariance}

We have now met three things that are secretly the same: the correlation
function $\xi$, the covariance matrix $C$, and the power spectrum $P$.
This short section makes the identification explicit, because it is the
conceptual hinge of the whole subject and it is what makes field
generation (\S\ref{sec:generate}) a two-line recipe.

\subsection{The covariance matrix and its diagonalisation}
\label{sec:cov-diag}

Recall from \S\ref{sec:xi-covariance} that the discretised field has
covariance matrix $C_{ij} = \xi(\bx_i - \bx_j)$, and that homogeneity
makes $C$ circulant (entries depend only on the offset $i - j$, wrapped
periodically).  The central fact about circulant matrices is:

\begin{quote}
\emph{Every circulant matrix is diagonalised by the discrete Fourier
transform, and its eigenvalues are the Fourier transform of its first
row.}
\end{quote}

The Fourier transform of the first row of $C$ --- that is, the Fourier
transform of $\xi$ --- is exactly $P(k)$ by Wiener--Khinchin
(\ref{eq:wiener-khinchin}).  So the power spectrum values are precisely
the \emph{eigenvalues of the covariance matrix}, and the Fourier modes
are its \emph{eigenvectors}:
\begin{equation}
  C\;\hat{\mathbf e}_\bk = P(\bk)\;\hat{\mathbf e}_\bk,
  \qquad
  (\hat{\mathbf e}_\bk)_j \propto e^{i\bk\cdot\bx_j}.
\end{equation}
This is Result~1 of \S\ref{sec:pk-derive} (modes are uncorrelated) stated
as linear algebra: in the Fourier basis the covariance matrix is
diagonal, with the power spectrum down the diagonal.

\subsection{Why this is the whole game}

The correlation function is the covariance in the position basis (a full
matrix); the power spectrum is the covariance in the Fourier basis (a
diagonal matrix).  ``Same operator, two bases.''  Everything that is
awkward about $\xi$ --- it couples every point to every other point ---
becomes trivial in the $P$ basis, where the modes decouple.  That is why:

\begin{itemize}
  \item \textbf{Theory} is usually written in $k$-space: the linear
        evolution equations, mode by mode, are ordinary differential
        equations with no coupling, because each Fourier mode evolves
        independently at linear order.
  \item \textbf{Field generation} is done in $k$-space: to build a field
        with a given covariance you take the square root of the
        covariance operator, and a diagonal matrix's square root is just
        the elementwise square root $\sqrt{P(\bk)}$
        (\S\ref{sec:generate}).
  \item \textbf{Measurement} of $P$ is a mode-by-mode variance
        (\S\ref{sec:estimate-pk}), needing no matrix inversion.
\end{itemize}

The rest of the note is really just consequences of this one
diagonalisation.

%======================================================================
\section{Gaussian random fields}
\label{sec:grf}

So far we have used only the first two moments (mean and covariance).
In general a random field has infinitely many independent statistics ---
the three-point function, four-point function, and so on --- and the
mean and covariance do not determine it.  \emph{Gaussian} random fields
are the special, enormously important case where they \emph{do}.  This
section explains why cosmological initial conditions are Gaussian to
excellent approximation, what ``Gaussian field'' means precisely, and
the theorem (Wick's) that makes the two-point function the whole story.

\subsection{Why the initial conditions are Gaussian}
\label{sec:grf-why}

The one-point distribution of $\delta$ --- the histogram of density
contrasts over all points --- is observed and predicted to be very
close to a Gaussian (a bell curve) in the early universe.  There are two
complementary reasons.

First, a \emph{central-limit} argument.  The density at a point is the
sum of contributions from a great many independent Fourier modes
(\S\ref{sec:fseries}).  The central limit theorem says that a sum of many
independent random contributions, none dominating, tends to a Gaussian
regardless of the distribution of the individual pieces.  So even if each
mode's amplitude were drawn from something non-Gaussian, their sum ---
the field at a point --- is pushed toward Gaussian.

Second, the \emph{physical origin}.  In the simplest models of inflation
the primordial fluctuations are quantum fluctuations of a nearly free
field, and a free field's ground-state fluctuations are exactly
Gaussian.  Deviations (``primordial non-Gaussianity'', parametrised by
$f_{\rm NL}$) are predicted to be tiny and have not been detected.  For
the purpose of generating initial conditions we take the field to be
Gaussian; \S\ref{sec:grf-fail} says a word about what changes if it is
not.

\subsection{What ``Gaussian random field'' means}
\label{sec:grf-def}

A zero-mean real field $\delta(\bx)$ is a \emph{Gaussian random field}
(GRF) if any one of the following equivalent statements holds.  They are
worth stating together because different ones are easiest to use in
different derivations.

\paragraph{Definition 1 (values at a finite set of points).}
Pick any finite set of points $\bx_1, \dots, \bx_n$.  The $n$ field
values $\bigl(\delta(\bx_1), \dots, \delta(\bx_n)\bigr)$ are jointly
drawn from a multivariate Gaussian distribution,
\begin{equation}
  p(\bu) = \frac{1}{\sqrt{(2\pi)^n \det C}}\,
           \exp\!\Bigl[-\tfrac12\,\bu^\top C^{-1}\,\bu\Bigr],
  \qquad C_{ij} = \xi(\bx_i - \bx_j),
  \label{eq:mvg}
\end{equation}
where $\bu = (u_1,\dots,u_n)^\top$ are the dummy values and $C$ is the
covariance matrix of Eq.~(\ref{eq:covmatrix}).  Note what appears in the
density: \emph{only} the covariance $C$.  The distribution has no other
parameters.  This is the precise sense in which ``mean and covariance
are everything'' --- for a Gaussian the density is literally built out
of $C$ and nothing else.

\paragraph{Definition 2 (every smoothing gives a Gaussian number).}
For every deterministic weight function $f(\bx)$, the single number
obtained by smoothing the field against it,
$F = \int f(\bx)\,\delta(\bx)\,\mathrm d^d x$, is an ordinary
one-dimensional Gaussian random variable.  (Definition~1 is the special
case where $f$ is a sum of spikes at the chosen points.)  This version is
convenient because it turns statements about the field into statements
about scalar Gaussians.

\paragraph{Definition 3 (the generating functional is Gaussian).}
Define the object that packages all the correlators at once,
\begin{equation}
  Z[J] \equiv \avg{\exp\!\Bigl(i\!\int J(\bx)\,\delta(\bx)\,\mathrm d^d x\Bigr)}.
  \label{eq:Z}
\end{equation}
For a GRF this equals
\begin{equation}
  Z[J] = \exp\!\Bigl[-\tfrac12\!\int\!\!\int J(\bx)\,\xi(\bx-\by)\,J(\by)\,
          \mathrm d^d x\,\mathrm d^d y\Bigr].
  \label{eq:Z-gauss}
\end{equation}
This is the field-theory generalisation of the elementary fact that a
one-dimensional zero-mean Gaussian of variance $\sigma^2$ has
characteristic function $\avg{e^{iJ\delta}} = e^{-\frac12 J^2\sigma^2}$.
The exponent is \emph{quadratic} in the source $J$, and everything below
follows from that.

\subsection{Wick's theorem: two points really are everything}
\label{sec:wick}

Definition~3 has a spectacular consequence.  The correlators of the
field are obtained by expanding $Z[J]$ in powers of $J$: the coefficient
of $i^n J(\bx_1)\cdots J(\bx_n)$ (times combinatorial factors) is the
$n$-point function $\avg{\delta(\bx_1)\cdots\delta(\bx_n)}$.  But the
right-hand side of Eq.~(\ref{eq:Z-gauss}) is the exponential of something
quadratic in $J$.  Expand it:
\begin{equation}
  Z[J] = 1 - \tfrac12\!\int\!\!\int J\,\xi\,J
         + \tfrac{1}{2!}\Bigl(\!-\tfrac12\!\int\!\!\int J\,\xi\,J\Bigr)^2
         - \cdots
\end{equation}
Every term is built entirely out of the two-point function $\xi$.
Matching powers of $J$ gives \emph{Wick's theorem}:
\begin{align}
  &\avg{\delta_1\,\delta_2\cdots\delta_{2k+1}} = 0
     &&\text{(odd number of fields),}
     \label{eq:wick-odd}\\[4pt]
  &\avg{\delta_1\,\delta_2\cdots\delta_{2k}}
     = \sum_{\text{pairings}}\;\prod_{\text{pairs }(i,j)} \xi(\bx_i - \bx_j)
     &&\text{(even number of fields),}
     \label{eq:wick-even}
\end{align}
where the sum runs over all ways of partitioning the fields into pairs.
The odd correlators vanish because a symmetric bell curve has no
lopsidedness (odd powers of $J$ never appear on the right of
Eq.~\ref{eq:Z-gauss}, whose exponent is even).  The even ones are sums
of products of two-point functions.

\paragraph{Extracting the four-point function explicitly.}  To see the
pairing structure actually emerge from Eq.~(\ref{eq:Z-gauss}), match the
$J^4$ terms on the two sides.  Abbreviate positions by integer labels and
write $\int\!\!\int J\,\xi\,J = \sum_{a,b}\xi_{ab}J_aJ_b$.  The definition
$Z[J] = \sum_n \frac{i^n}{n!}\sum \avg{\delta_1\cdots\delta_n}
J_1\cdots J_n$ has $n=4$ term
$\frac{1}{4!}\sum_{abcd}\avg{\delta_a\delta_b\delta_c\delta_d}
J_aJ_bJ_cJ_d$.  On the Gaussian side, the only source of four powers of
$J$ is the second term of the expanded exponential,
\begin{equation}
  \frac{1}{2!}\Bigl(-\tfrac12\sum_{a,b}\xi_{ab}J_aJ_b\Bigr)^{\!2}
  = \frac{1}{8}\sum_{a,b,c,d}\xi_{ab}\,\xi_{cd}\,J_aJ_bJ_cJ_d .
\end{equation}
The product $J_aJ_bJ_cJ_d$ is symmetric under permuting the four labels,
so only the symmetrised coefficient survives.  Of the $4! = 24$
relabelings of $(a,b,c,d)$, exactly $8$ reproduce each of the three
distinct pairings --- $2$ from swapping within the first pair, $2$ within
the second, $2$ from exchanging the two pairs (using
$\xi_{ab}=\xi_{ba}$) --- so
$\sum_{\text{perms}}\xi_{ab}\xi_{cd}
= 8(\xi_{ab}\xi_{cd}+\xi_{ac}\xi_{bd}+\xi_{ad}\xi_{bc})$.  Equating the
symmetric $J^4$ coefficients, $\frac{1}{4!}\avg{\delta_a\delta_b\delta_c
\delta_d} = \frac{1}{8}\cdot\frac{8}{24}
(\xi_{ab}\xi_{cd}+\cdots)$, the factors of $24$ cancel and
\begin{equation}
  \avg{\delta_1\delta_2\delta_3\delta_4}
  = \xi_{12}\,\xi_{34} + \xi_{13}\,\xi_{24} + \xi_{14}\,\xi_{23},
  \label{eq:wick4}
\end{equation}
the three ways of pairing four objects.  The same bookkeeping at order
$J^{2k}$ gives Eq.~(\ref{eq:wick-even}) with its
$(2k)!/(2^k k!)$ pairings --- $15$ for the six-point function, and so on.

\paragraph{The punchline.}  Every correlator of a Gaussian field is an
explicit polynomial in $\xi$.  Nothing is independent of the two-point
function.  Combined with Definition~3 (which says $Z[J]$, hence the
entire probability distribution, is fixed once $\xi$ is known), this
means:
\begin{quote}
\textbf{A Gaussian random field is completely specified by its mean and
its two-point function (equivalently, its power spectrum).}
\end{quote}
This is why cosmologists spend so much effort measuring $P(k)$: for
Gaussian initial conditions it is not just \emph{a} statistic, it is
\emph{the} statistic --- knowing it is knowing the field's full
distribution.  It is also what makes generating a realisation so cheap
(\S\ref{sec:generate}): match the two-point function and, for free, you
have matched everything.

\subsection{Where Gaussianity fails}
\label{sec:grf-fail}

If the primordial field is not Gaussian --- for instance the local-model
form
\begin{equation}
  \delta_{\rm NG}(\bx) = \delta_{\rm G}(\bx)
       + f_{\rm NL}\bigl[\delta_{\rm G}^2(\bx) - \avg{\delta_{\rm G}^2}\bigr],
  \label{eq:fnl}
\end{equation}
with $\delta_{\rm G}$ Gaussian --- then Eq.~(\ref{eq:Z-gauss}) no longer
holds, the higher cumulants are nonzero, and $P(k)$ no longer determines
the distribution.  The three-point function (bispectrum) becomes an
independent observable, and a single ``draw white noise, colour by
$\sqrt P$'' construction cannot produce the field.  Everything in the
next section assumes Gaussianity; non-Gaussian generators bolt an extra
nonlinear term like Eq.~(\ref{eq:fnl}) onto the Gaussian core.

%======================================================================
\section{Generating a Gaussian random field on a computer}
\label{sec:generate}

We now turn the theory into an algorithm.  The goal: produce a realisation
$\delta(\bx)$ on a grid whose power spectrum is a prescribed $P(k)$.
Because the field is Gaussian, matching $P(k)$ matches the full
distribution (\S\ref{sec:wick}), so this single recipe is a complete
construction of the initial conditions.

\subsection{The idea: take the square root of the covariance}
\label{sec:gen-idea}

Start in the finite-dimensional picture of \S\ref{sec:xi-covariance}: we
want a Gaussian vector $\boldsymbol\delta \in \mathbb R^N$ with a
prescribed covariance $C = \avg{\boldsymbol\delta\,\boldsymbol\delta^\top}$.
The trick is to start from the easiest possible Gaussian vector ---
\emph{white noise} $\boldsymbol\mu$, with independent unit-variance
components, $\avg{\boldsymbol\mu\,\boldsymbol\mu^\top} = I$ --- and
transform it linearly.  If $C = LL^\top$ for some matrix $L$ (a ``square
root'' of $C$), then
\begin{equation}
  \boldsymbol\delta = L\,\boldsymbol\mu
  \quad\Longrightarrow\quad
  \avg{\boldsymbol\delta\,\boldsymbol\delta^\top}
     = L\,\avg{\boldsymbol\mu\,\boldsymbol\mu^\top}\,L^\top
     = L\,I\,L^\top = LL^\top = C.
  \label{eq:sqrt-cov}
\end{equation}
A linear map of a Gaussian is Gaussian (Definition~2), so
$\boldsymbol\delta$ is a Gaussian vector with exactly the covariance we
wanted.  The whole problem reduces to factoring $C$.

For a general covariance the factorisation $L$ (Cholesky, or the
symmetric square root from the eigendecomposition) costs $\mathcal
O(N^3)$ --- hopeless for a field with $N = 256^3 \approx 10^7$ cells.
But our $C$ is not general: it is circulant (\S\ref{sec:cov-diag}).  And
we already know a circulant matrix's eigenbasis --- the Fourier basis ---
and its eigenvalues --- the power spectrum.  Its square root is therefore
trivial: in Fourier space $C$ is the diagonal matrix
$\mathrm{diag}\,P(\bk)$, and its square root is
$\mathrm{diag}\,\sqrt{P(\bk)}$.  Applying $L$ is then just:
Fourier-transform, multiply each mode by $\sqrt{P(\bk)}$, transform back.
Cost: $\mathcal O(N \log N)$, the cost of an FFT.  This is the payoff of
the diagonalisation of \S\ref{sec:covariance}.

\subsection{White noise and its flat spectrum}
\label{sec:whitenoise}

Unit-variance white noise $\mu(\bq)$ is the field of independent
identically distributed $\mathcal N(0,1)$ values, one per cell, with
\begin{equation}
  \avg{\mu(\bq)} = 0,
  \qquad
  \avg{\mu(\bq)\,\mu(\bq')} = \delta^D(\bq - \bq')
  \;\;\text{(continuum)},
  \quad
  \avg{\mu_\bn\,\mu_{\bn'}} = \delta^K_{\bn\bn'}
  \;\;\text{(grid)}.
  \label{eq:whitenoise}
\end{equation}
Its correlation function is a spike at zero separation and zero
elsewhere: each cell knows nothing about its neighbours.  By
Wiener--Khinchin (\ref{eq:wiener-khinchin}) the Fourier transform of a
delta is a constant, so white noise has a \emph{flat} power spectrum,
$P_\mu(k) = 1$ --- equal power on every scale.  That is the origin of the
name (white light has equal power at every frequency).  White noise is
the canonical Gaussian field: the identity covariance, the thing we
colour.

A fact we will need for normalisation: the (unnormalised) discrete
Fourier transform of grid white noise has a known per-mode variance.
With $\tilde\mu_\bk = \sum_\bn \mu_\bn e^{-i\bk\cdot\bx_\bn}$ over
$N_{\rm tot}$ cells,
\begin{align}
  \avg{|\tilde\mu_\bk|^2}
  &= \sum_{\bn,\bn'} \avg{\mu_\bn\,\mu_{\bn'}}\,
     e^{-i\bk\cdot(\bx_\bn - \bx_{\bn'})}
   = \sum_{\bn,\bn'} \delta^K_{\bn\bn'}\,
     e^{-i\bk\cdot(\bx_\bn - \bx_{\bn'})}
   = \sum_\bn 1 = N_{\rm tot}.
  \label{eq:whitenoise-permode}
\end{align}
The Kronecker delta collapses the double sum to a single one; the result
is independent of $\bk$, which is the Fourier-space statement of
``white''.  This $N_{\rm tot}$ is the factor that has to be divided out
when converting a raw FFT amplitude into a physically normalised power
spectrum; \code{cosmo\_ic.tex} Appendix~A
track exactly where it goes.

\subsection{Colouring: the recipe}
\label{sec:colouring}

Putting \S\ref{sec:gen-idea} and \S\ref{sec:whitenoise} together gives
the standard algorithm.  Two equivalent versions, differing only in
whether the randomness is drawn in real space or Fourier space.

\paragraph{Method A --- white noise in real space, then colour.}
\begin{enumerate}
  \item Draw a real white-noise field: $\mu_\bn \sim \mathcal N(0,1)$,
        independently in every cell.
  \item Forward-FFT to get $\tilde\mu_\bk$.
  \item Colour: multiply each mode by $\sqrt{P(\bk)\,V}$, giving
        $\tilde\delta_\bk = \sqrt{P(\bk)\,V}\,\tilde\mu_\bk$.  (Zero the
        $\bk = 0$ mode to enforce exactly zero mean.)
  \item Inverse-FFT; take the real part.
\end{enumerate}
Because the input $\mu$ was real, its transform is automatically
Hermitian-symmetric (\S\ref{sec:hermitian}), and multiplying by the
real-valued $\sqrt{P(\bk)\,V}$ preserves that symmetry, so the inverse
transform is real with no hand-tuning.  This automatic reality is the
practical advantage of Method~A, and it is why simulation codes use it.
Checking the per-mode variance with
Eq.~(\ref{eq:whitenoise-permode}) confirms the normalisation:
$\avg{|\tilde\delta_\bk|^2} = P(\bk)\,V\,\avg{|\tilde\mu_\bk|^2}
= P(\bk)\,V\,N_{\rm tot}$, which combines with the inverse-FFT and
cell-size factors to reproduce Eq.~(\ref{eq:pk-var}); the careful factor
accounting is Appendix~A of \code{cosmo\_ic.tex}.

\paragraph{Method B --- draw the modes directly in Fourier space.}
\begin{enumerate}
  \item For each independent mode $\bk$, draw a complex Gaussian
        amplitude $\tilde\delta_\bk = \sqrt{P(\bk)\,V/2}\,(u + iv)$ with
        $u, v \sim \mathcal N(0,1)$ independent, so that
        $\avg{|\tilde\delta_\bk|^2} = P(\bk)\,V$.
  \item Impose Hermitian symmetry by hand:
        $\tilde\delta_{-\bk} = \tilde\delta_\bk^*$.
  \item For the self-conjugate modes ($\bk = 0$ and the Nyquist corners
        on an even grid) draw a purely real amplitude.
  \item Inverse-FFT to real space.
\end{enumerate}
Method~B needs more bookkeeping but lets you set individual modes by hand
--- useful for ``fixed-and-paired'' variance-reduction
\citep{AnguloPontzen2016} and for zoom initial conditions --- and it
saves one FFT.  On a large box the two methods are statistically
identical; they differ only in their treatment of the largest modes on a
small box, where Method~A convolves $P$ with the box window and Method~B
does not.  \code{ic\_sampling.tex} and \citet{Pen1997, Sirko2005} discuss
that ``$P$-sampled versus $\xi$-sampled'' distinction, which can bias
$\sigma_8$ in small boxes.

\subsection{The statistics of a single Fourier amplitude}
\label{sec:mode-stats}

It is worth knowing the distribution --- not just the mean --- of a
single mode's amplitude, because it is what one plots to check that a
generated field is correctly normalised, and it explains the ragged
look of a raw (unbinned) power-spectrum measurement.

Fix a generic mode $\bk$ (not $\bk=0$, not a Nyquist self-conjugate
corner).  From Method~B, $\tilde\delta_\bk = X + iY$ with $X$ and $Y$
independent zero-mean Gaussians of equal variance $\sigma_k^2 =
P(\bk)V/2$.  Two classic results follow purely from the geometry of a
two-dimensional Gaussian.

\paragraph{The amplitude is Rayleigh-distributed.}  Write
$\tilde\delta_\bk$ as a point $(X, Y)$ in the complex plane, drawn from
the isotropic 2D Gaussian
\begin{equation}
  p(X, Y)\,\mathrm dX\,\mathrm dY
    = \frac{1}{2\pi\sigma_k^2}
      \exp\!\Bigl(-\frac{X^2+Y^2}{2\sigma_k^2}\Bigr)\,\mathrm dX\,\mathrm dY.
\end{equation}
Change to polar coordinates $R = |\tilde\delta_\bk|$, $\theta =
\arg\tilde\delta_\bk$, with $\mathrm dX\,\mathrm dY = R\,\mathrm
dR\,\mathrm d\theta$.  The phase $\theta$ comes out uniform on
$[0, 2\pi)$ and independent of $R$ (isotropy), and marginalising over it
(the $\int_0^{2\pi}\mathrm d\theta = 2\pi$ cancels the $1/2\pi$
prefactor) leaves
\begin{equation}
  p(R) = \frac{R}{\sigma_k^2}\,
         \exp\!\Bigl(-\frac{R^2}{2\sigma_k^2}\Bigr), \qquad R \ge 0,
  \label{eq:rayleigh}
\end{equation}
the \emph{Rayleigh distribution}.  The phase of each mode is completely
random --- this is why a Gaussian field is sometimes said to have
``random phases'' --- and the amplitude is Rayleigh.

\paragraph{The power in a mode is exponentially distributed.}  Change
variable once more to the mode power $U = R^2 = |\tilde\delta_\bk|^2$,
with $\mathrm dU = 2R\,\mathrm dR$.  The factor $R$ in
Eq.~(\ref{eq:rayleigh}) is exactly cancelled by the Jacobian, leaving a
pure exponential:
\begin{equation}
  p(U) = \frac{1}{2\sigma_k^2}\,
         \exp\!\Bigl(-\frac{U}{2\sigma_k^2}\Bigr), \qquad U \ge 0,
  \label{eq:exponential}
\end{equation}
with mean $\avg{U} = 2\sigma_k^2 = P(\bk)V$, consistent with
Eq.~(\ref{eq:pk-var}).  The exponential has a heavy tail: the most likely
mode power is zero, the median is $\ln 2 \approx 0.69$ times the mean,
and rare modes carry several times the mean.  A single realisation's
$|\tilde\delta_\bk|^2$ therefore scatters wildly around $P(\bk)$ --- with
$100\%$ relative fluctuation per mode --- which is why a power-spectrum
estimate must average many modes to be useful (\S\ref{sec:estimate-pk}).
Figure~\ref{fig:grf-exp} shows the three distributions --- the circular
cloud of $(X, Y)$, the Rayleigh $R$, and the exponential $R^2$ ---
measured from the FFT of real white noise.

\begin{figure}[t]
\centering
\includegraphics[width=\textwidth]{grf_exponential.pdf}
\caption{Per-mode statistics of the FFT of real white noise on a
$64\times64$ grid, over $20{,}000$ realisations at one generic mode
(reproduced from a companion note).
\textbf{(a)} The complex amplitude $(X, Y) = (\mathrm{Re},\mathrm{Im})$
is an isotropic 2D Gaussian --- random phase, Gaussian real and imaginary
parts.
\textbf{(b)} The modulus $R = |\tilde\delta_\bk|$ follows the Rayleigh
distribution of Eq.~(\ref{eq:rayleigh}).
\textbf{(c)} The mode power $R^2 = |\tilde\delta_\bk|^2$ is exponentially
distributed (Eq.~\ref{eq:exponential}; note the log $y$-axis).  Red
curves are the analytic predictions.  The wide exponential scatter is why
a power-spectrum estimator must bin many modes.}
\label{fig:grf-exp}
\end{figure}

\subsection{Worked example: a Gaussian-correlated field in 1D}
\label{sec:grf-example}

To make the recipe concrete, build a one-dimensional field with a chosen
correlation length $\lambda$ and check by hand that every property comes
out right.  Take the target correlation function to be a Gaussian bump,
\begin{equation}
  \xi(r) = A\,e^{-r^2/\lambda^2},
  \label{eq:ex-xi}
\end{equation}
so that $A = \xi(0) = \avg{\delta^2}$ is the variance and $\lambda$ is
the $1/e$ correlation length.  Its power spectrum is the Fourier
transform (Wiener--Khinchin), a standard Gaussian integral:
\begin{equation}
  P(k) = \int_{-\infty}^{\infty}\!\mathrm dr\; A\,e^{-r^2/\lambda^2}\,
         e^{-ikr}
       = A\,\lambda\sqrt\pi\;e^{-k^2\lambda^2/4}.
  \label{eq:ex-Pk}
\end{equation}
(A short check: transforming Eq.~(\ref{eq:ex-Pk}) back with the inverse
of Eq.~(\ref{eq:P-xi-pair}) returns Eq.~(\ref{eq:ex-xi}), using
$\int\mathrm dk\,e^{-ak^2}e^{ikr} = \sqrt{\pi/a}\,e^{-r^2/4a}$ with
$a = \lambda^2/4$.)  A short-correlation-length field (small $\lambda$)
has a broad power spectrum (power out to large $k$); a long-correlation
field has power concentrated at small $k$ --- the reciprocal relationship
between real-space width and Fourier-space width.

\paragraph{Discrete recipe (Method A).}  On a periodic grid of $N$ cells
of spacing $\Delta$:
\begin{enumerate}
  \item Draw $\mu_n \sim \mathcal N(0,1)$, $n = 0,\dots,N-1$.
  \item FFT: $\tilde\mu_m = \sum_n \mu_n e^{-2\pi i mn/N}$.
  \item Colour mode by mode: $\tilde\delta_m =
        \sqrt{P(k_m)/(N\Delta)}\,\tilde\mu_m$, with $k_m = 2\pi m/(N\Delta)$.
  \item Inverse-FFT: $\delta_n = N^{-1}\sum_m \tilde\delta_m e^{+2\pi i mn/N}$.
\end{enumerate}
Three properties can be verified by hand on a small grid (try $N = 8$,
$\lambda = 2\Delta$):
\begin{itemize}
  \item[(a)] \textbf{Each $\delta_n$ is Gaussian}, because step~4 makes it
        a fixed linear combination of the Gaussian $\mu$'s.
  \item[(b)] \textbf{The variance is right}: summing
        $\avg{|\tilde\delta_m|^2} = P(k_m)/\Delta$ over $m$ and dividing
        by the inverse-FFT $N^2$ gives $\avg{\delta_n^2} = (N\Delta)^{-1}
        \sum_m P(k_m) \to \int\mathrm dk\,P(k)/(2\pi) = A$ in the
        continuum limit --- exactly $\xi(0)$.
  \item[(c)] \textbf{The two-point function is right}: the same sum with
        a phase gives $\avg{\delta_n\delta_{n+r}} \to A\,e^{-r^2\Delta^2/
        \lambda^2} = \xi(r\Delta)$.
\end{itemize}
By Wick (\S\ref{sec:wick}), once (a) and (c) hold the \emph{entire}
distribution of $\{\delta_n\}$ matches the target.  Matching the two
lines of second-moment algebra has, for free, matched every higher
correlator.  That is the whole content of Gaussianity, made arithmetic.

%======================================================================
\section{Measuring $\xi$ and $P$ from one realisation}
\label{sec:estimate}

The last piece is the inverse problem: given one box (a simulation
snapshot or a galaxy survey), estimate its correlation function and power
spectrum.  This is where ergodicity (\S\ref{sec:ergodic}) does its work
--- we replace the unavailable ensemble average with an average over the
box --- and where the finite size of the box sets the ultimate error.
The estimator details specific to $N$-body simulations (mass assignment,
window deconvolution, shot noise) live in Appendix~B of
\code{cosmo\_ic.tex}; here we cover the ideas.

\subsection{Estimating the power spectrum}
\label{sec:estimate-pk}

Equation~(\ref{eq:pk-var}) says $P(\bk) = \avg{|\hat\delta_\bk|^2}/V$.
The ensemble average is unavailable, but under isotropy all modes with
the same $|\bk|$ have the same $P$, so we \emph{average over the modes in
a shell} $k \le |\bk| < k + \Delta k$ as a stand-in for the ensemble
average:
\begin{equation}
  \hat P(k) = \frac{1}{V\,N_k}\sum_{\bk\,\in\,\text{shell}} |\hat\delta_\bk|^2,
  \label{eq:pk-estimator}
\end{equation}
where $N_k$ is the number of modes in the shell.  Because each
$|\hat\delta_\bk|^2$ is exponentially distributed
(\S\ref{sec:mode-stats}) with $100\%$ scatter, averaging $N_k$ of them
cuts the fractional error to $\sim 1/\sqrt{N_k}$.  Large-scale (small-$k$)
shells contain few modes and are noisy; small-scale (large-$k$) shells
contain many modes and are smooth.  The procedure in practice: assign the
particles to a density grid, FFT, square, and bin in $|\bk|$.

\subsection{Cosmic variance: the price of a finite box}
\label{sec:cosmicvar}

The number of independent modes in a shell of width $\Delta k$ is
\begin{equation}
  N_k \approx \frac{4\pi k^2\,\Delta k}{(2\pi/L)^3}
            = \frac{V k^2\,\Delta k}{2\pi^2},
\end{equation}
because $k$-space is quantised in cells of size $(2\pi/L)^3$.  Larger box
$\to$ finer $k$-spacing $\to$ more modes per shell $\to$ smaller error.
On the largest scales (smallest $k$) there are only a handful of modes no
matter how well you simulate, and the fractional error $\sim
1/\sqrt{N_k}$ is irreducible.  This is \emph{cosmic variance}: the
fundamental uncertainty from having sampled a finite piece of the random
field.  It is the box-scale echo of the per-mode exponential scatter of
\S\ref{sec:mode-stats}, and it is why large-scale measurements (in
surveys and in simulations alike) are the hardest to pin down.

\subsection{Estimating the correlation function}
\label{sec:estimate-xi}

Two routes, using the two faces of the same object.

\paragraph{Direct pair counting.}  Straight from the definition
(\ref{eq:xi-pair}): count the pairs of particles separated by $r$
(within a bin), and compare to the number expected for a random
distribution,
\begin{equation}
  \hat\xi(r) = \frac{DD(r)}{RR(r)} - 1,
\end{equation}
where $DD$ counts data--data pairs and $RR$ the pairs of a uniform random
catalogue with the same geometry.  (Real surveys use the
Landy--Szalay combination $(DD - 2DR + RR)/RR$, which cancels
edge effects to leading order.)  Pair counting is $\mathcal O(N^2)$ but
gives $\xi(r)$ directly at the separations of interest.  It is the right
tool at low redshift where the clustering signal is large; at high
redshift the signal in $\xi$ is minuscule and the power spectrum (or the
grid estimator) is the better probe.

\paragraph{Via the FFT.}  Because $\xi$ and $P$ are a Fourier pair
(\ref{eq:P-xi-pair}), a second route is to measure $P(k)$ by
Eq.~(\ref{eq:pk-estimator}) and inverse-transform it.  This is $\mathcal
O(N\log N)$ and gives $\xi$ at all separations at once, but it inherits
the grid and window effects of the FFT.  Which route wins depends on the
scales of interest and the number density; the two are consistent where
they overlap.

\subsection{Grid assignment and its window}
\label{sec:cic}

To FFT a set of particles you must first paint them onto a grid.  The
common ``cloud-in-cell'' (CIC) scheme spreads each particle over its
eight nearest cells, which is a convolution of the true field with a
small assignment kernel.  Convolution in real space is multiplication in
Fourier space, so the measured $\hat\delta_\bk$ is the true one times a
\emph{window function} $W(\bk)$ that must be divided back out
(``deconvolved'') before reading off $P(k)$.  The assignment also
introduces \emph{shot noise} --- a flat $1/\bar n$ floor from the
discreteness of the particles --- that must be subtracted, and
\emph{aliasing} of power from above the grid Nyquist frequency.  These
three corrections (window, shot noise, aliasing) are exactly the subject
of Appendix~B of \code{cosmo\_ic.tex}; they are the difference between a
raw FFT and a publishable power spectrum.  One cosmology-specific subtlety
worth flagging: lattice initial conditions have \emph{sub-Poissonian}
shot noise (the particles start on a regular grid, not at random
positions), so the naive $1/\bar n$ subtraction is wrong at early times
--- see the same appendix.

%======================================================================
\section{Summary and where to go next}
\label{sec:summary}

The logical spine of this note, in one paragraph.  A cosmological density
field is a \emph{random field}: predictable only in its statistics, of
which --- under homogeneity, isotropy, and (for measurement) ergodicity
--- the leading one is the two-point function.  In real space that is the
\emph{correlation function} $\xi(r)$, the excess clustering of pairs and
equivalently the covariance of the field between two points.  In Fourier
space, homogeneity makes the modes uncorrelated and the same information
becomes the \emph{power spectrum} $P(k)$, the variance per mode; the two
are a Fourier pair (Wiener--Khinchin) and are literally the covariance
operator written in two bases.  For a \emph{Gaussian} field --- which
inflation delivers and the central limit theorem reinforces --- Wick's
theorem makes this two-point function the \emph{entire} story: mean plus
$P(k)$ fixes the whole distribution.  That is what lets us both
\emph{generate} a realisation (draw white noise, colour it by
$\sqrt{P(k)}$ with an FFT --- the square root of the diagonalised
covariance) and \emph{measure} $P(k)$ back from a single box (bin the
squared Fourier amplitudes), with the finite box size setting an
irreducible cosmic-variance floor.

\medskip\noindent Where each thread continues in the other notes:
\begin{itemize}
  \item \textbf{The cosmological content} --- how $P(k)$ arises from
        inflation and the transfer function, how the field is displaced
        into particle positions (Zel'dovich, 2LPT), the physical
        correlation and velocity statistics $\xi(r)$ and $\psi(r)$:
        \code{cosmo\_ic.tex}.  Its Appendix~A does the careful Fourier
        factor bookkeeping quoted here; its Appendix~B is the estimator
        detail of \S\ref{sec:estimate}.
  \item \textbf{The numerical Fourier transform} --- the discrete and
        fast transforms, per-mode statistics, radix and parallel FFT
        algorithms: any standard FFT reference.  Figure~\ref{fig:grf-exp} is borrowed
        from there.
  \item \textbf{Sampling methods and finite-box biases} --- the
        $P$-sampled versus $\xi$-sampled distinction of
        \S\ref{sec:colouring}, box-window truncation: \code{ic\_sampling.tex}.
  \item \textbf{The generator internals} --- how MUSIC2 realises the
        $\delta = T\star\mu$ construction of \S\ref{sec:generate} on a
        refinement hierarchy: \code{music2\_internals.tex}.
\end{itemize}

%======================================================================
\section*{Further reading}
\addcontentsline{toc}{section}{Further reading}

Ordered roughly by increasing formality; a good first path is
\citet{MoBoschWhite2010} \S4.3--4.4, then \citet{BBKS1986} Appendix~A,
then \citet{Pen1997}.

\begin{itemize}
  \item \citet{MoBoschWhite2010}, \emph{Galaxy Formation and Evolution},
        \S4.3--4.4: the cleanest textbook chapter on cosmological
        Gaussian random fields --- definitions, correlation function,
        power spectrum, peak statistics --- at exactly the level of this
        note.
  \item \citet{Peacock1999}, \emph{Cosmological Physics}, ch.~16: the
        central-limit origin of the Gaussian one-point distribution and
        the spectral $\delta = T\star\mu$ construction.
  \item \citet{DodelsonSchmidt2020}, \emph{Modern Cosmology}, ch.~6--7:
        connects the random-field formalism to the inflationary origin of
        $P(k)$ and the transfer function.
  \item \citet{BBKS1986} (Bardeen, Bond, Kaiser \& Szalay), Appendix~A:
        a self-contained physicist's primer on multivariate Gaussians,
        $n$-point functions, and Wick contractions.
  \item \citet{vanKampen2007}, \emph{Stochastic Processes in Physics and
        Chemistry}, ch.~3: the clearest account of Gaussian processes,
        the characteristic functional (Definition~3), and why all
        cumulants beyond the second vanish.
  \item \citet{Pen1997}: the original FFT recipe for generating a
        cosmological Gaussian random field --- the construction of
        \S\ref{sec:generate}.  \citet{Sirko2005} analyses the finite-box
        sampling biases of \S\ref{sec:colouring}.
  \item \citet{AnguloPontzen2016}: the ``fixed-and-paired'' construction
        that exploits the mode phase/amplitude freedom to suppress cosmic
        variance.
\end{itemize}

%======================================================================
\begin{thebibliography}{99}

\bibitem[{Angulo~\& Pontzen(2016)}]{AnguloPontzen2016}
Angulo, R.~E.\ \& Pontzen, A. 2016, MNRAS, 462, L1.

\bibitem[{Bardeen et al.(1986)}]{BBKS1986}
Bardeen, J.~M., Bond, J.~R., Kaiser, N., \& Szalay, A.~S. 1986,
ApJ, 304, 15.

\bibitem[{Dodelson~\& Schmidt(2020)}]{DodelsonSchmidt2020}
Dodelson, S.\ \& Schmidt, F. 2020, ``Modern Cosmology'' 2nd ed.\
(Academic Press).

\bibitem[{Mo et al.(2010)}]{MoBoschWhite2010}
Mo, H., van den Bosch, F., \& White, S. 2010, ``Galaxy Formation
and Evolution'' (Cambridge University Press).

\bibitem[{Peacock(1999)}]{Peacock1999}
Peacock, J.~A. 1999, ``Cosmological Physics'' (Cambridge
University Press).

\bibitem[{Pen(1997)}]{Pen1997}
Pen, U.-L. 1997, ApJS, 109, 1.

\bibitem[{Sirko(2005)}]{Sirko2005}
Sirko, E. 2005, ApJ, 634, 728.

\bibitem[{van~Kampen(2007)}]{vanKampen2007}
van Kampen, N.~G. 2007, ``Stochastic Processes in Physics and
Chemistry'' 3rd ed.\ (North-Holland).

\end{thebibliography}

\end{document}
